\documentclass[11pt,a4paper]{article}
\usepackage[T1]{fontenc}
\usepackage{lmodern,microtype}
\usepackage[margin=27mm,headheight=14pt]{geometry}
\usepackage{amsmath,amssymb,amsthm,mathtools,bm}
\usepackage{booktabs,longtable,array,enumitem}
\usepackage[table]{xcolor}
\usepackage{fancyhdr}
\usepackage{hyperref}
\hypersetup{colorlinks=true,linkcolor=blue!45!black,citecolor=blue!45!black,
 urlcolor=blue!45!black,pdftitle={When Dougall Tails Terminate},
 pdfauthor={AI-assisted research continuation prepared for Vladimir Reshetnikov}}
\usepackage{bookmark}
\numberwithin{equation}{section}
\newtheorem{theorem}{Theorem}[section]
\newtheorem{proposition}[theorem]{Proposition}
\newtheorem{lemma}[theorem]{Lemma}
\newtheorem{corollary}[theorem]{Corollary}
\theoremstyle{definition}\newtheorem{definition}[theorem]{Definition}
\theoremstyle{remark}\newtheorem{remark}[theorem]{Remark}
\DeclareMathOperator{\Li}{Li}
\DeclareMathOperator{\arctanh}{arctanh}
\DeclareMathOperator{\Res}{Res}
\newcommand{\C}{\mathbb C}\newcommand{\Q}{\mathbb Q}\newcommand{\Z}{\mathbb Z}
\newcommand{\R}{\mathbb R}\newcommand{\N}{\mathbb N}
\newcommand{\dd}{\,\mathrm d}
\newcommand{\Gcal}{\mathcal G}\newcommand{\Rcal}{\mathcal R}
\newcommand{\Bcal}{\mathcal B}\newcommand{\Lcal}{\mathfrak L}
\newcommand{\Repart}{\operatorname{Re}}
\newcommand{\pin}{\texttt{da2792fc9cc8ef19c6a314be6c65e1437fbdc66b}}
\newcommand{\coef}[1]{\left[#1\right]}
\allowdisplaybreaks[2]
\emergencystretch=2em
\pagestyle{fancy}\fancyhf{}
\fancyhead[L]{\small When Dougall tails terminate}
\fancyhead[R]{\small ProveIt continuation}
\fancyfoot[C]{\thepage}
\begin{document}
\begin{titlepage}
\thispagestyle{empty}
\vspace*{17mm}
{\Large\sffamily PROVEIT / EXACT SPECIAL-FUNCTION IDENTITIES\par}
\vspace{13mm}
{\huge\bfseries When Dougall Tails Terminate\par}
\vspace{6mm}
{\Large A rationality classification, Stieltjes cancellation,\par
and an all-depth harmonic--polylogarithmic sum formula\par}
\vspace{13mm}
{\large Research continuation prepared for Vladimir Reshetnikov\par}
\vspace{2mm}
{\large AI-assisted mathematical draft\par}
\vspace{2mm}
{11 October 2026\par}
\vspace{12mm}
\begin{abstract}
We classify all rational and all finitely terminating centered Gamma
quotients at the integral resonances of the positive-parameter
Rogers--Dougall family. Rationality occurs precisely when one of three
symmetric parameters is integral. Finite termination occurs precisely on
the evident short-integer polynomial loci or on one exceptional symmetric
Laurent-polynomial locus. A rational difference-equation argument shows
why termination of the formal asymptotic series is an exact identity,
not merely an asymptotic assertion. In contrast, analytic curves of
nonrational summands can have identically zero summed remainder.

A two-parameter resolvent specialization then separates the Stieltjes
boundary term from the convergent higher moments. Subtracting two
resolvents cancels every Stieltjes coefficient and yields a Gamma
expression for all fixed-order polygamma moments. At the half-center,
a normalized Gamma quotient evaluates an infinite family of weighted
composition sums in ordinary generalized harmonic numbers. The resulting
all-depth level-two polylogarithmic sum formula has values in the algebra
of ordinary zeta values and $\log 2$. Explicit cubic and quartic sums,
a dilogarithmic integral, Euler primitives, and zero-safe derivative
formulas are given with proofs and reproducible diagnostics.

The classical summation and the repository's prior all-resonance
continuation are credited separately. Global literature priority is not
established, and no proof-assistant formalization is claimed.
\end{abstract}
\vfill
{\small Repository snapshot: \pin\par}
{\small This package is additive. No remote repository files were modified.\par}
\end{titlepage}
\tableofcontents
\newpage

\section{Scope, antecedents, and the precise additions}
\label{sec:scope}
The ProveIt manuscript organizes polylogarithms, harmonic sums, Gamma
products, and Hurwitz-zeta derivatives by exact transformations rather
than by numerical recognition alone. Its canonical mixed-Dougall section
proves a first-resonance identity, fixed mixed derivatives, and a weighted
trigamma square \cite{repo-mixed}. A preserved continuation already proves
the centered all-resonance subtraction, its collapsed digamma value, and
all spectral jets \cite{repo-centered}. Those results are inputs and
antecedents here, not newly solved conjectures.

The present work begins with a different question: \emph{when is the
subtracted tail actually zero term by term?} Several sufficient polynomial
loci appear in the preserved report. We prove a complete classification
within its positive-parameter domain. We then show that a vanishing
\emph{sum} of the tail is a substantially weaker condition: there are
analytic families of nonrational Gamma quotients whose remainders sum
to zero. This gives a concrete distinction between an exact finite
summand identity and an exact infinite cancellation.

The second addition is a resolvent construction. A symmetric perturbation
of the Gamma parameters produces the denominator
$x_n(x_n^2-v^2)$. Its Gamma evaluation is an ordinary convergent
identity. In a companion boundary identity, every generalized Stieltjes
constant can be displayed explicitly. Taking a resolvent difference
removes all those constants at once. After half-center normalization,
this yields a compact generating theorem for nonlinear harmonic sums,
together with exact signed sums of multiple polylogarithms of arbitrary
depth.

The Rogers--Dougall summation is classical \cite{dlmf16}, as are Gamma
recurrences, Stirling expansions, and Hurwitz-zeta identities
\cite{dlmf5,dlmf25}. Parameter differentiation of hypergeometric sums
is an established source of harmonic and zeta identities. The Euler-sum
literature includes extensive linear and quadratic evaluations and
colored multiple-zeta representations; Xu and Wang provide one relevant
modern reference \cite{xu-wang}. Consequently, the statements below are
proved mathematical results, but ``proved here'' does not mean
``first in the literature.'' The rationality/termination classification,
the nonterminating zero-remainder curves, and the specific unified
resolvent/composition organization are the principal additions relative
to the inspected repository text.

\paragraph{Audit limits.}
The source inspection covered the pinned canonical text and the specific
preserved continuation listed in the source-audit file. The five ZIP
archives present in \texttt{docs/incoming} were inventoried by name and
metadata but could not be unpacked through the available binary-access
route. Their contents were not checked for overlap. The external
antecedent check used DLMF and the Xu--Wang publisher page, not an
exhaustive literature review. The Gaussian $S_6$ conjecture is not
resolved by this package. The manuscript's rejected frozen $S_8$ vector
is not revived. These limitations do not weaken the proofs below, but
they do limit priority and coverage claims.

\section{Centered Dougall subtraction: the established starting point}
\label{sec:starting}
Throughout the classification sections, fix
\begin{equation}
 \kappa,b,c>0,\qquad d_0=1+2\kappa-b-c>0,\qquad x_n=n+\kappa.
 \label{eq:positive-domain}
\end{equation}
All denominator Gamma factors mean entire reciprocal-Gamma functions.
In particular, a zero reciprocal-Gamma value must not be replaced by a
zero germ before differentiation. Define
\begin{align}
 W_n(u)={}&x_n\frac{\Gamma(n+2\kappa)\Gamma(n+b)\Gamma(n+c)
                         \Gamma(n+d_0-u)}
 {\Gamma(n+1)\Gamma(n+1+2\kappa-b)\Gamma(n+1+2\kappa-c)
                         \Gamma(n+b+c+u)},\label{eq:W}\\
 P(u)={}&\frac{\Gamma(b)\Gamma(c)\Gamma(d_0-u)}
 {\Gamma(d_0)\Gamma(b+u)\Gamma(c+u)}.\label{eq:P}
\end{align}
The classical Rogers--Dougall identity in this normalization is
\begin{equation}
 \sum_{n\ge0}W_n(u)=\frac12\Gamma(u)P(u),
 \qquad 0<\Repart u<d_0.
 \label{eq:dougall}
\end{equation}
The hypergeometric excess is $2u$. The factor $1/2$ is essential: the
usual very-well-poised factor is $(n+\kappa)/\kappa$, whereas
\eqref{eq:W} uses $x_n$ and unnormalized Gamma factors.

Put
\begin{equation}
 \alpha(u)=(\kappa,b-\kappa,c-\kappa,1+\kappa-b-c-u).
 \label{eq:alpha}
\end{equation}
The denominator shifts relative to $x_n$ are $1-\alpha_i(u)$.
Let $B_j(x)$ denote Bernoulli polynomials and define
\begin{align}
 \ell_j(u)&=-\frac{\sum_{i=1}^4B_{2j+1}(\alpha_i(u))}{j(2j+1)},
 \qquad j\ge1,\label{eq:ell}\\
 A_0(u)&=1,\qquad
 A_j(u)=\frac1j\sum_{h=1}^j h\ell_h(u)A_{j-h}(u).\label{eq:A}
\end{align}
Shifted Stirling expansion and Bernoulli reflection give
\begin{equation}
 W_n(u)=x_n^{-1-2u}
 \left(\sum_{j=0}^{K-1}A_j(u)x_n^{-2j}+O(x_n^{-2K})\right).
 \label{eq:centered-expansion}
\end{equation}
For fixed $K$, the remainder is locally uniform in pole-free compact
parameter sets. Every fixed parameter derivative is controlled on a
slightly larger compact set by Cauchy's formula. Only even inverse powers
occur in the correction series.

\begin{proposition}[Prior centered subtraction and collapse]
\label{prop:prior}
For $K\ge1$ and $-K<\Repart u<d_0$,
\begin{align}
 \Rcal_K(u)&=\sum_{n\ge0}\left[
 W_n(u)-\sum_{j=0}^{K-1}A_j(u)x_n^{-1-2u-2j}\right]\notag\\
 &=\frac12\Gamma(u)P(u)-\sum_{j=0}^{K-1}
                   A_j(u)\zeta(1+2u+2j,\kappa).
 \label{eq:subtraction}
\end{align}
The bracketed series and all its fixed derivatives converge locally
normally. If $N\ge0$, then
\begin{align}
 A_N(-N)=\rho_N
 &=\frac{(-1)^N}{N!}(d_0)_N(1-b)_N(1-c)_N,\label{eq:rho}\\
 \Rcal_{N+1}(-N)
 &=\frac{\rho_N}{2}\bigl[\psi(N+1)-\psi(d_0+N)
              -\psi(b)-\psi(c)+2\psi(\kappa)\bigr].\label{eq:collapse}
\end{align}
Here $(a)_N$ is the rising factorial. These formulas include $\rho_N=0$.
\end{proposition}
\begin{proof}[Proof, recalled to specify the input]
The compact-uniform remainder in \eqref{eq:centered-expansion} is
summable on the stated strip. On $\Repart u>0$, the sum follows by
\eqref{eq:dougall}; analytic continuation gives
\eqref{eq:subtraction}. Its holomorphic left side proves removal of
all apparent poles on the right.

At $u=-N$, compare residues. The Gamma term has residue $\rho_N/2$;
the only polar zeta term is $j=N$, with residue $A_N(-N)/2$.
This proves \eqref{eq:rho}, first for generic parameters and then
polynomially. The cancellation of the finite parts is encoded by
\begin{equation}
 (\partial_u-D)A_N(u)=\sum_{h=1}^N\frac{B_{2h}(\kappa)}h A_{N-h}(u),
 \qquad D=\partial_\kappa+\partial_b+\partial_c.
 \label{eq:coefficient-law}
\end{equation}
Indeed $Dd_0=0$, and $\partial_u-D$ acts on the four entries of
\eqref{eq:alpha} by $(-1,0,0,0)$. Apply this operator to
$\exp(\sum_{j\ge1}\ell_j y^j)=\sum_{j\ge0}A_jy^j$.
The Bernoulli derivative identity gives \eqref{eq:coefficient-law}.

The constant term of the Gamma expression is
\[
 \frac{\rho_N}{2}
 [\psi(N+1)-\psi(d_0+N)-\psi(b-N)-\psi(c-N)].
\]
The polar zeta counterterm contributes
$A'_N(-N)/2-\rho_N\psi(\kappa)$. The lower zeta values are
Bernoulli polynomials; \eqref{eq:coefficient-law} cancels them and
replaces the coefficient derivative by $D\rho_N$. Finally,
\[
 D\rho_N=\rho_N\sum_{h=1}^N
              \left(\frac1{b-h}+\frac1{c-h}\right),
\]
and the digamma recurrence produces \eqref{eq:collapse}.
Continuity proves the exceptional cases. This is the centered
continuation's proof mechanism \cite{repo-centered}.
\end{proof}

\section{A complete rationality and finite-termination classification}
\label{sec:classification}
At a fixed integral resonance, put
\begin{equation}
 d=d_0+N,\qquad
 \alpha=(\kappa,b-\kappa,c-\kappa,d-\kappa),\qquad
 b+c+d=1+2\kappa+N.
 \label{eq:resonance-parameters}
\end{equation}
The meromorphic function whose sampled values are $W_n(-N)$ is
\begin{equation}
 F_N(x)=x\prod_{i=1}^4
                  \frac{\Gamma(x+\alpha_i)}{\Gamma(x+1-\alpha_i)}.
 \label{eq:F}
\end{equation}
Here ``rational'' means belonging to $\mathbb R(x)$; it does not
require rational coefficients or rational values at integer arguments.
The three parameters $b,c,d$ are symmetric in this expression, although
$d_0>0$ singles out one of them in the original convergence domain.

\subsection{Two general lemmas}
\begin{lemma}[Rational Gamma quotients]
\label{lem:gamma-rational}
For real shifts $a_1,\ldots,a_r$ and $e_1,\ldots,e_r$, the quotient
$\prod_i\Gamma(x+a_i)/\Gamma(x+e_i)$ is rational in $x$ if and only
if the two multisets of shifts agree modulo $\Z$.
\end{lemma}
\begin{proof}
If the multisets agree, pair equal residue classes and use
$\Gamma(z+h)/\Gamma(z)=(z)_h$ for integral $h$, interpreting negative
$h$ by the reciprocal product. This gives a rational function.
Conversely, in each residue class of shifts modulo $\Z$, numerator
Gamma factors contribute one pole at every sufficiently far negative
point of that arithmetic progression. Denominator Gamma factors
contribute the corresponding zeros. Unless the multiplicities agree,
the quotient has infinitely many zeros or poles in that progression.
A nonzero rational function cannot do so. Finite initial cancellations
cannot change this conclusion.
\end{proof}

\begin{lemma}[Formal termination is exact]
\label{lem:formal-exact}
Suppose a meromorphic function $F$ has a nonzero finite Laurent
asymptotic expansion $P$ to all algebraic orders on the positive real
axis, and $F(x+1)=R(x)F(x)$ for a rational function $R$. Then $F=P$,
provided $F(x)/P(x)\to1$ as $x\to+\infty$.
\end{lemma}
\begin{proof}
For every prescribed power $M$, the difference $F(x)-P(x)$ is
$O(x^{-M})$. Substitution into the difference equation implies that
$P(x+1)-R(x)P(x)$ decays faster than every power. This is a rational
function; therefore it vanishes identically. Away from the finite
exceptional sets, $F/P$ is consequently $1$-periodic. For each
sufficiently large real $x$,
\[
 \frac{F(x)}{P(x)}=\frac{F(x+j)}{P(x+j)}\longrightarrow1
 \qquad(j\to\infty).
\]
Thus the quotient equals one on an interval, and meromorphic
continuation proves the identity. This argument is what rules out
an unexamined beyond-all-orders correction.
\end{proof}

For \eqref{eq:F}, the required rational difference equation is
\begin{equation}
 \frac{F_N(x+1)}{F_N(x)}
 =\frac{x+1}{x}\prod_{i=1}^4\frac{x+\alpha_i}{x+1-\alpha_i}.
 \label{eq:difference}
\end{equation}
Its leading asymptotic power is $x^{2N-1}$.

\subsection{The classification theorem}
\begin{theorem}[Rationality and termination]
\label{thm:classification}
Assume \eqref{eq:positive-domain} and $N\ge0$.

\textup{(i)} The function $F_N(x)$ is rational if and only if at least
one of $b,c,d$ is an integer.

\textup{(ii)} Its formal centered inverse-power expansion terminates
if and only if $F_N$ is a Laurent polynomial. This occurs precisely
in either of the following cases:
\begin{equation}
 \text{some member of }\{b,c,d\}\text{ lies in }\{1,\ldots,N\};
 \label{eq:short-integer}
\end{equation}
or, up to permutation,
\begin{equation}
 (b,c,d)=(N+1,\kappa,\kappa).
 \label{eq:exceptional}
\end{equation}
In the first case $F_N$ is a polynomial. In the second case,
\begin{equation}
 F_N(x)=\frac1x\prod_{j=1}^{N}\bigl(x^2-(\kappa-j)^2\bigr).
 \label{eq:exceptional-product}
\end{equation}
The empty product at $N=0$ is one. Overlap of the two cases is allowed.
\end{theorem}
\begin{proof}
Since $\sum_i\alpha_i=N+1$, Lemma~\ref{lem:gamma-rational} says that
rationality is equivalent to invariance of the four-element multiset
$\{\alpha_i\bmod\Z\}$ under negation. Such an invariant multiset,
with integral total sum, has a partner for $\alpha_1$ among the
other three elements, with integral pair sum. To see the only delicate
case, suppose $\alpha_1$ is a fixed class under negation, hence $0$
or $1/2$. If no second copy of that fixed class occurs, the other
nonfixed classes pair off, leaving an odd number of the other fixed
class. The total sum is then half-integral, a contradiction.
The possible partner sums are exactly $b,c,d$.

Conversely, if one of those three pair sums is integral, the
complementary pair sum is integral as well, because their total is
$N+1$. Pairing proves negation invariance and hence rationality.
This establishes (i).

If the formal expansion terminates, it is a finite Laurent polynomial
with leading coefficient one. Equation~\eqref{eq:difference} and
Lemma~\ref{lem:formal-exact} identify it exactly with $F_N$.
The converse is immediate from uniqueness of the asymptotic expansion.
It remains to classify the rational functions which are Laurent
polynomials.

Choose an integer member $m$ of $\{b,c,d\}$ and rename the other two
$c_*,d_*$. Positivity implies $m\ge1$. If $m\le N$, Gamma recurrence
gives
\begin{equation}
 F_N(x)=x\prod_{j=1}^{m-1}\bigl(x^2-(\kappa-j)^2\bigr)
          \prod_{r=1}^{N-m}\bigl(x^2-(\kappa-c_*+r)^2\bigr).
 \label{eq:polynomial-product}
\end{equation}
This is the short-integer polynomial locus.

Suppose no integer member is at most $N$, and choose $m$ to be the
\emph{smallest integer member}. Put $L=m-N\ge1$. The same recurrence
now gives the rational expression
\begin{equation}
 F_N(x)=x\,
 \frac{\displaystyle\prod_{j=1}^{m-1}(x^2-(\kappa-j)^2)}
 {\displaystyle\prod_{r=0}^{L-1}(x^2-(\kappa-c_*-r)^2)}.
 \label{eq:rational-product}
\end{equation}
A numerator and denominator factor can cancel only if
\[
 c_*+r=j\quad\text{or}\quad d_*+L-1-r=j,
 \qquad 1\le j\le m-1,
\]
where the second equality uses $c_*+d_*=1+2\kappa-L$.
If the parameter on the left is nonintegral, neither equality is
possible. If integral, it is at least $m$ by minimality of $m$,
so neither equality is possible again. Thus no quadratic factors
cancel.

A Laurent polynomial can have no pole away from $x=0$. Consequently,
every denominator root in \eqref{eq:rational-product} must be zero.
The consecutive numbers $\kappa-c_*-r$, $0\le r<L$, can all vanish
only when $L=1$ and $c_*=\kappa$. The parameter sum then forces
$d_*=\kappa$ and $m=N+1$. Formula~\eqref{eq:exceptional-product}
follows. This proves necessity and sufficiency, including multiple
roots and zero factors by meromorphic continuation.
\end{proof}

\begin{remark}[Exact scope of the symmetric description]
Under $d_0>0$, one has $d>N$. Thus the short-integer alternative is
actually attained by $b$ or $c$ in the original parameter chart.
The exceptional locus $d=N+1$, $b=c=\kappa$, $d_0=1$ is admissible
for every $\kappa>0$. The permutations with $b=N+1$ or $c=N+1$
require $\kappa>N$. The symmetric theorem is useful, but the
positivity domain must still be enforced after a permutation.
\end{remark}

\subsection{Consequences and examples}
A terminating centered expression need not be a polynomial: the generic
exceptional expression has a simple $x^{-1}$ term. Its residue coefficient
is
\begin{equation}
 \rho_N=(-1)^N\prod_{j=1}^N(\kappa-j)^2
 \quad\text{on }(b,c,d)=(\kappa,\kappa,N+1).
 \label{eq:rho-exceptional}
\end{equation}
This agrees with \eqref{eq:rho}. On all terminating loci,
\begin{equation}
 W_n(-N)=\sum_{j=0}^N A_j(-N)x_n^{2N-1-2j},\qquad
 \Rcal_{N+1}(-N)=0,
 \label{eq:termwise-zero}
\end{equation}
with any unused final coefficients zero.

For example, $N=2$, $\kappa=b=c=1/2$, $d_0=1$ gives
\[
 F_2(x)=\frac{(x^2-1/4)(x^2-9/4)}x
       =x^3-\frac52x+\frac9{16x}.
\]
By contrast, $N=1$, $\kappa=2$, $b=3$, $c=3/4$, $d_0=5/4$
is rational but does not terminate:
\[
 F_1(x)=\frac{x^3(x^2-1)}{(x^2-25/16)(x^2-1/16)}.
\]
Both statements are identities of meromorphic functions; evaluating only
at large integer indices is not needed for the classification.

\section{Zero summed tails without rational summands}
\label{sec:zero-curves}
The implication from termwise termination to a zero sum has no converse.
The collapsed expression in Proposition~\ref{prop:prior} permits a precise
local construction rather than a numerical search for examples.

\begin{theorem}[Nonrational zero-remainder curves]
\label{thm:zero-curves}
For every $N\ge0$ and $0<\kappa<1$, there is a real analytic even
function $s(v)$ near $v=0$, with $s(0)=0$, such that the parameters
\begin{equation}
 b=\kappa+s(v)+v,\quad c=\kappa+s(v)-v,\quad
 d_0=1-2s(v)
 \label{eq:zero-curve}
\end{equation}
satisfy
\begin{equation}
 \Rcal_{N+1}(-N)=0.
 \label{eq:zero-curve-value}
\end{equation}
For all sufficiently small real $v\ne0$, the corresponding $F_N$ is
not rational and its centered expansion does not terminate. Moreover,
\begin{equation}
 s''(0)=\frac{\psi''(\kappa)}{\psi'(N+1)-\psi'(\kappa)}>0.
 \label{eq:zero-curve-curvature}
\end{equation}
\end{theorem}
\begin{proof}
Consider the real analytic function
\begin{align*}
 B(s,v)={}&\psi(N+1)-\psi(N+1-2s)\\
          &-\psi(\kappa+s+v)-\psi(\kappa+s-v)+2\psi(\kappa).
\end{align*}
At the origin, $B=0$,
\[
 B_s=2[\psi'(N+1)-\psi'(\kappa)]\ne0,
 \qquad B_v=0,\qquad B_{vv}=-2\psi''(\kappa).
\]
The trigamma is strictly decreasing on the positive axis, since
$\psi''(x)=-2\sum_{n\ge0}(n+x)^{-3}<0$. The analytic implicit
function theorem gives a unique analytic $s(v)$ solving $B(s(v),v)=0$.
Evenness follows from uniqueness and $B(s,-v)=B(s,v)$.
Differentiating twice gives \eqref{eq:zero-curve-curvature}; in
particular $s(v)>0$ for small nonzero real $v$.

All parameters remain positive. Shrink the neighborhood so that
$b,c\in(0,1)$ and $N<d=N+1-2s(v)<N+1$ for $v\ne0$.
No member of $\{b,c,d\}$ is then integral, and $\rho_N\ne0$.
Equation~\eqref{eq:collapse} proves \eqref{eq:zero-curve-value},
while Theorem~\ref{thm:classification} proves nonrationality and
nontermination. The argument uses an exact implicit family, not the
accuracy of a numerical zero.
\end{proof}

This theorem is not a correction to an asserted converse in the inspected
report: no such converse was found. It is a new separation result and an
audit guard. A small or zero total residual cannot be used to infer
summand-level collapse.

\section{A symmetric resolvent and its Stieltjes boundary}
\label{sec:resolvent}
We now use $a>0$ for the center and set $\kappa=a$, $b=a+v$,
$c=a-v$, $d_0=1$. Let $|v|<a$ and define
\begin{align}
 T_n(u;a)&=\frac{\Gamma(n+2a)\Gamma(n+1-u)}
                    {\Gamma(n+1)\Gamma(n+2a+u)},\label{eq:Tn}\\
 Q_a(u,v)&=\frac{\pi u}{\sin\pi u}
 \frac{\Gamma(a+v)\Gamma(a-v)}
      {\Gamma(a+v+u)\Gamma(a-v+u)}.\label{eq:Qa}
\end{align}
Then $W_n(u)=x_nT_n(u;a)/(x_n^2-v^2)$, where $x_n=n+a$.
The symmetry makes the resolvent denominator appear without any
approximation.

\begin{theorem}[Ordinary convergent resolvent identity]
\label{thm:resolvent}
For $a>0$, $|v|<a$, and $-1<\Repart u<1$,
\begin{equation}
 \sum_{n\ge0}\frac{v^2T_n(u;a)}{x_n(x_n^2-v^2)}
       =\frac{Q_a(u,v)-Q_a(u,0)}{2u}.
 \label{eq:resolvent}
\end{equation}
The right side is interpreted removably at $u=0$. The series and
all fixed $u$- and $v$-derivatives converge locally normally.
\end{theorem}
\begin{proof}
In $0<\Repart u<1$, subtract two instances of
\eqref{eq:dougall}, one with $v$ and one with zero. The elementary
identity
\[
 \frac{x}{x^2-v^2}-\frac1x=\frac{v^2}{x(x^2-v^2)}
\]
and $\Gamma(1+u)\Gamma(1-u)=\pi u/\sin\pi u$ give
\eqref{eq:resolvent}.

The centered expansion gives
$T_n(u;a)=x_n^{-2u}(1+O(x_n^{-2}))$ locally uniformly.
Thus the summand is $O(n^{-3-2\Repart u})$ on compact sets of the
stated domain. Cauchy's formula on slightly larger compact sets gives
normal convergence of each fixed derivative. The factors
$1/\Gamma(n+2a+u)$ are entire and introduce no poles; the numerator
Gamma factors are pole-free for $\Repart u<1$. Both sides therefore
continue to the full strip. Since $Q_a(0,v)=1$, the apparent pole
at zero is removable.
\end{proof}

The subtraction suppresses the only divergent moment. To see exactly
what is removed, retain that moment and use
\begin{equation}
 \zeta(1+s,a)=\frac1s+\sum_{m\ge0}\frac{(-1)^m\gamma_m(a)}{m!}s^m.
 \label{eq:stieltjes-convention}
\end{equation}

\begin{theorem}[All Stieltjes boundary jets]
\label{thm:boundary}
On the same domain as Theorem~\ref{thm:resolvent},
\begin{align}
 \Bcal_a(u,v)&=\sum_{n\ge0}
 \left[\frac{x_nT_n(u;a)}{x_n^2-v^2}-x_n^{-1-2u}\right]\notag\\
 &=\frac{Q_a(u,v)}{2u}-\zeta(1+2u,a).
 \label{eq:boundary}
\end{align}
The bracketed series is locally normally convergent, and for every
integer $m\ge0$,
\begin{equation}
 \partial_u^m\Bcal_a(0,v)
 =\frac{\partial_u^{m+1}Q_a(0,v)}{2(m+1)}-(-2)^m\gamma_m(a).
 \label{eq:boundary-jets}
\end{equation}
Subtracting the $v=0$ identity cancels $\gamma_m(a)$ at every order.
\end{theorem}
\begin{proof}
The leading centered term of the bracket is canceled, leaving the same
summable order as in Theorem~\ref{thm:resolvent}. The identity follows
first in the convergent strip and then by continuation. Expand both
terms on the right at zero. The common coefficient $1/(2u)$ cancels;
the coefficient of $u^m$, multiplied by $m!$, is
\eqref{eq:boundary-jets}. Termwise differentiation was established
before the expansion.
\end{proof}

Write $\psi_\pm=\psi(a\pm v)$ and
$\psi'_\pm=\psi'(a\pm v)$. The first two boundary values are
\begin{align}
 \Bcal_a(0,v)&=\psi(a)-\frac{\psi_++\psi_-}{2},\label{eq:boundary0}\\
 \partial_u\Bcal_a(0,v)&=
 \frac{(\psi_++\psi_-)^2-\psi'_+-\psi'_-+2\zeta(2)}4
                  +2\gamma_1(a).\label{eq:boundary1}
\end{align}
At $v=0$, the second formula becomes the ordinary convergent identity
\begin{equation}
 \sum_{n\ge0}\frac{2\log(n+a)-\psi(n+1)-\psi(n+2a)}{n+a}
 =\psi(a)^2-\frac12\psi'(a)+\frac12\zeta(2)+2\gamma_1(a).
 \label{eq:stieltjes-linear}
\end{equation}
For complex $a$ in a pole-free continuation, the square is algebraic,
not an absolute square. The numerical implementation in this package
uses positive real $a$ for Stieltjes constants.

\section{Finite polygamma formulas for every higher resolvent moment}
\label{sec:polygamma}
For $k\ge1$, expansion of \eqref{eq:resolvent} in $v$ gives
\begin{equation}
 \sum_{n\ge0}\frac{T_n(u;a)}{(n+a)^{2k+1}}
 =\frac{[v^{2k}]Q_a(u,v)}{2u}.
 \label{eq:moment-generator}
\end{equation}
This is already a finite Gamma-derivative expression. An explicit
coefficient algebra makes its all-orders content transparent. Put
\begin{equation}
 D_j(u)=\frac{2}{(2j)!}
             [\psi^{(2j-1)}(a)-\psi^{(2j-1)}(a+u)].
 \label{eq:Dj}
\end{equation}
Taylor expansion in $v$ yields
\begin{equation}
 \frac{[v^{2k}]Q_a(u,v)}{2u}
 =\frac{Q_a(u,0)}{2u}[z^k]
       \exp\left(\sum_{j\ge1}D_j(u)z^j\right).
 \label{eq:finite-bell}
\end{equation}
Only $j\le k$ are needed. The numerator vanishes at $u=0$ before
division, because every $D_j(0)$ is zero.

Let $Y_m$ be the complete exponential Bell polynomial, defined by
\[
 \exp\left(\sum_{r\ge1}g_r\frac{u^r}{r!}\right)
       =\sum_{m\ge0}Y_m(g_1,\ldots,g_m)\frac{u^m}{m!}.
\]
For the left side of \eqref{eq:moment-generator}, the Bell data are
\begin{equation}
 g_{n,r}=(-1)^r\psi^{(r-1)}(n+1)-\psi^{(r-1)}(n+2a).
 \label{eq:gnr}
\end{equation}
Consequently every sum
\begin{equation}
 \sum_{n\ge0}\frac{Y_m(g_{n,1},\ldots,g_{n,m})}{(n+a)^{2k+1}},
 \qquad m\ge0,\ k\ge1,
 \label{eq:bell-moment}
\end{equation}
is given by the $m$th derivative of the finite expression
\eqref{eq:finite-bell}. All its constants are polygamma values at $a$
and ordinary even zeta values. This does not assert a reduction of
arbitrary unrelated polygamma products; it evaluates this particular
Bell family.

\begin{corollary}[An all-weight linear identity]
\label{cor:linear}
For $a>0$ and every integer $k\ge1$,
\begin{align}
 &\sum_{n\ge0}\frac{\psi(n+1)+\psi(n+2a)}{(n+a)^{2k+1}}
 \notag\\
 &\quad=\frac{2k+1}{2}\zeta(2k+2,a)
        +2\psi(a)\zeta(2k+1,a)
        -\sum_{j=1}^{k-1}\zeta(2j+1,a)\zeta(2k-2j+1,a).
 \label{eq:linear-general}
\end{align}
\end{corollary}
\begin{proof}
The required sum is the negative first $u$-derivative of
\eqref{eq:moment-generator}, since
$\partial_uT_n(0;a)=-\psi(n+1)-\psi(n+2a)$.
In \eqref{eq:finite-bell}, retain terms through degree two in $u$:
\[
 D_j(u)=-\frac{2\psi^{(2j)}(a)}{(2j)!}u
       -\frac{\psi^{(2j+1)}(a)}{(2j)!}u^2+O(u^3),
 \quad Q_a(u,0)=1-2\psi(a)u+O(u^2).
\]
The exponential contributes a linear term and the quadratic convolution
of two $D_j$'s. Insert
$\psi^{(r)}(a)=(-1)^{r+1}r!\zeta(r+1,a)$ for $r\ge1$.
Coefficient extraction gives \eqref{eq:linear-general}.
\end{proof}

\begin{corollary}[A quadratic identity at every positive center]
\label{cor:quadratic}
For $a>0$,
\begin{align}
 &\sum_{n\ge0}\frac{[\psi(n+1)+\psi(n+2a)]^2
                    +\psi'(n+1)-\psi'(n+2a)}{(n+a)^3}\notag\\
 &\quad=-\frac16\psi^{(4)}(a)+\psi(a)\psi^{(3)}(a)
       -[2\psi(a)^2+\zeta(2)-\psi'(a)]\psi''(a).
 \label{eq:quadratic-general}
\end{align}
\end{corollary}
\begin{proof}
For $k=1$, the right side of \eqref{eq:finite-bell} is
$Q_a(u,0)[\psi'(a)-\psi'(a+u)]/(2u)$. Its second derivative at
zero is the displayed expression, using
\[
 Q_a(u,0)=1-2\psi(a)u+
 [2\psi(a)^2+\zeta(2)-\psi'(a)]u^2+O(u^3).
\]
The second derivative of $T_n$ is exactly the numerator in the
summand. Normal convergence justifies the operation.
\end{proof}

\section{A Gamma generator for nonlinear harmonic sums}
\label{sec:harmonic}
The half-center $a=1/2$ removes the shifted polygamma arguments after
one normalization. Define
\begin{equation}
 P_n(u)=\frac{(1-u)_n}{(1+u)_n}
       =\prod_{j=1}^n\frac{j-u}{j+u},\qquad P_0(u)=1,
 \label{eq:Pn}
\end{equation}
and
\begin{equation}
 \Gcal(u,t)=\frac{\Gamma(1+u)^2
                   \Gamma((1+t)/2)\Gamma((1-t)/2)}
 {\Gamma((1+t)/2+u)\Gamma((1-t)/2+u)}.
 \label{eq:G}
\end{equation}

\begin{theorem}[Normalized half-center generator]
\label{thm:half-generator}
For $|t|<1$ and $\Repart u>-1$,
\begin{equation}
 \sum_{n\ge0}\frac{t^2P_n(u)}{(2n+1)((2n+1)^2-t^2)}
       =\frac{\Gcal(u,t)-\Gcal(u,0)}{4u}.
 \label{eq:half-resolvent}
\end{equation}
The singularity at $u=0$ is removable. In particular, for every $k\ge1$,
\begin{equation}
 \sum_{n\ge0}\frac{P_n(u)}{(2n+1)^{2k+1}}
       =\frac{[t^{2k}]\Gcal(u,t)}{4u}.
 \label{eq:half-moments}
\end{equation}
Every fixed Taylor coefficient at $u=0$ can be summed termwise.
\end{theorem}
\begin{proof}
Set $a=1/2$ and $v=t/2$ in \eqref{eq:resolvent}, and multiply by
$\Gamma(1+u)/\Gamma(1-u)$. Then $T_n$ becomes $P_n$ and $Q_a$
becomes $\Gcal$. The rational summand acquires a factor two when
$x_n=(2n+1)/2$ is substituted; this is the source of $4u$ rather
than $2u$ on the right of \eqref{eq:half-resolvent}.

Initially the calculation applies on $-1<\Repart u<1$ away from
removable factors. The finite product \eqref{eq:Pn} is holomorphic
for $\Repart u>-1$. On compact subsets, the Gamma-ratio form gives
$P_n(u)=O(n^{-2\Repart u})$, uniformly also across positive integral
$u$ by using the entire reciprocal-Gamma factor. The sum in
\eqref{eq:half-resolvent} is therefore locally normally convergent
on the larger half-plane. The Gamma expression is holomorphic there
apart from its removable $u=0$ singularity. Analytic continuation
extends the identity. Expansion in $t$ and Cauchy's formula prove
\eqref{eq:half-moments} and its coefficient assertions.
\end{proof}

\subsection{Weighted compositions of generalized harmonic numbers}
For a composition $\mathbf r=(r_1,\ldots,r_d)$, define
\begin{equation}
 H_n(\mathbf r)=\sum_{n\ge j_1>\cdots>j_d\ge1}
                   \frac1{j_1^{r_1}\cdots j_d^{r_d}},
 \qquad H_n(\varnothing)=1.
 \label{eq:strict-H}
\end{equation}
The notation $\mathbf r\vDash m$ means that all $r_i\ge1$ and
$\sum_i r_i=m$. For $m=0$ there is the single empty composition.
Put
\begin{equation}
 C_m(n)=\sum_{\mathbf r\vDash m}2^{\ell(\mathbf r)}H_n(\mathbf r).
 \label{eq:Cm}
\end{equation}

\begin{lemma}[Finite composition identity]
\label{lem:composition}
For every $n\ge0$,
\begin{equation}
 P_n(u)=\sum_{m\ge0}(-1)^m C_m(n)u^m
 \label{eq:P-composition}
\end{equation}
as a Taylor series near zero.
\end{lemma}
\begin{proof}
Each finite factor is
$(1-u/j)/(1+u/j)=1+2\sum_{r\ge1}(-1)^r(u/j)^r$.
Choose the nonconstant terms at a strictly decreasing collection of
indices and sum their positive exponents. The resulting coefficient
is exactly \eqref{eq:Cm}. There is no infinite summation over $n$
in this finite identity.
\end{proof}

For $H_n=H_n^{(1)}$ and $H_n^{(r)}=\sum_{j=1}^n j^{-r}$,
\begin{align}
 C_0(n)&=1,& C_1(n)&=2H_n,& C_2(n)&=2H_n^2,\notag\\
 C_3(n)&=\frac{4H_n^3+2H_n^{(3)}}3,&
 C_4(n)&=\frac{2H_n^4+4H_nH_n^{(3)}}3.
 \label{eq:low-C}
\end{align}
These also follow from
\begin{equation}
 \log P_n(u)=-2\sum_{j\ge0}\frac{H_n^{(2j+1)}}{2j+1}u^{2j+1}.
 \label{eq:log-Pn}
\end{equation}
Thus only odd-index harmonic numbers enter the logarithm, although
the equivalent composition expression runs over all compositions.

\begin{theorem}[All-weight harmonic composition sum formula]
\label{thm:harmonic}
For every $m\ge0$ and $k\ge1$,
\begin{equation}
 \boxed{\displaystyle
 \sum_{n\ge0}\frac{C_m(n)}{(2n+1)^{2k+1}}
       =\frac{(-1)^m}{4}[u^{m+1}t^{2k}]\Gcal(u,t).}
 \label{eq:all-harmonic}
\end{equation}
Every value belongs to $\Q[\log2,\zeta(2),\zeta(3),\ldots]$.
An explicit finite expression is obtained by exponentiating
\begin{align}
 \log\Gcal(u,t)={}&4u\log2
 +\sum_{r\ge2}\frac{2(-1)^r(2-2^r)}r\zeta(r)u^r\notag\\
 &+\sum_{\substack{r\ge1\\j\ge1}}
 \frac{2(-1)^{r+1}}r\binom{r+2j-1}{2j}
 (2^r-2^{-2j})\zeta(r+2j)u^rt^{2j}.
 \label{eq:log-G}
\end{align}
Only terms of $u$-degree at most $m+1$ and $t$-degree at most $2k$
are required. The resulting expression is homogeneous of weight
$m+2k+1$ when $\log2$ has weight one and $\zeta(r)$ has weight $r$.
\end{theorem}
\begin{proof}
Combine \eqref{eq:half-moments} and Lemma~\ref{lem:composition};
normal convergence permits coefficient extraction. Equivalently,
$C_m(n)$ is bounded by a constant times $(1+\log n)^m$ for fixed
$m$, so the displayed ordinary sum converges absolutely.

For completeness, expand the logarithm of \eqref{eq:G} in $u$.
Its coefficient of $u^r$ is
\[
 \frac{2\psi^{(r-1)}(1)
   -\psi^{(r-1)}((1+t)/2)-\psi^{(r-1)}((1-t)/2)}{r!}.
\]
The $r=1$, $t=0$ coefficient is $4\log2$. At $t=0$ and $r\ge2$,
use
$\psi^{(r-1)}(1)=(-1)^r(r-1)!\zeta(r)$ and
$\psi^{(r-1)}(1/2)=(-1)^r(r-1)!(2^r-1)\zeta(r)$.
For the coefficient of $t^{2j}$, expand the symmetric pair and use
the same half-argument formula at order $r+2j-1$. Simplification gives
\eqref{eq:log-G}.

Every monomial in the logarithm has constant weight equal to its
$u$-degree plus its $t$-degree. Exponentiation preserves this grading.
The coefficient in \eqref{eq:all-harmonic} therefore has the asserted
weight. These are formal grading statements; they do not assert
arithmetic independence of the constants.
\end{proof}

\subsection{Explicit linear, quadratic, cubic, and quartic identities}
Write $L=\log2$. The first four nontrivial rows of
\eqref{eq:all-harmonic}, at $k=1$, simplify to
\begin{align}
 \sum_{n\ge0}\frac{H_n}{(2n+1)^3}
 &=\frac{45}{32}\zeta(4)-\frac74L\zeta(3),\label{eq:example1}\\
 \sum_{n\ge0}\frac{H_n^2}{(2n+1)^3}
 &=\frac{31}{8}\zeta(5)-\frac{45}{8}L\zeta(4)
   +\frac72L^2\zeta(3)-\frac78\zeta(2)\zeta(3),\label{eq:example2}\\
 \sum_{n\ge0}\frac{2H_n^3+H_n^{(3)}}{(2n+1)^3}
 &=-14L^3\zeta(3)+\frac{21}{2}L\zeta(2)\zeta(3)
      -\frac{21}{4}\zeta(3)^2\notag\\
 &\quad+\frac{135}{4}L^2\zeta(4)-\frac{93}{2}L\zeta(5)
      +\frac{945}{64}\zeta(6),\label{eq:example3}\\
 \sum_{n\ge0}\frac{H_n^4+2H_nH_n^{(3)}}{(2n+1)^3}
 &=14L^4\zeta(3)-21L^2\zeta(2)\zeta(3)+21L\zeta(3)^2\notag\\
 &\quad-45L^3\zeta(4)+93L^2\zeta(5)-\frac{945}{16}L\zeta(6)\notag\\
 &\quad-\frac{39}{2}\zeta(3)\zeta(4)-\frac{93}{4}\zeta(2)\zeta(5)
      +\frac{1143}{16}\zeta(7).\label{eq:example4}
\end{align}
The proof of each is finite coefficient extraction from a proved
analytic identity. The exact script performs that extraction in a
polynomial algebra and checks these simplified presentations using only
the even-zeta relations. The displayed combinations, not $H_n^3$ and
$H_n^4$ separately, are what the theorem evaluates. Nor are these sums
the manuscript's alternating Gaussian sums $S_p$.

\section{All-depth polylogarithms, integrals, and primitives}
\label{sec:polylog}
We use decreasing-index multiple polylogarithms
\begin{equation}
 \Li_{r_1,\ldots,r_d}(z_1,\ldots,z_d)
 =\sum_{j_1>\cdots>j_d\ge1}
 \frac{z_1^{j_1}\cdots z_d^{j_d}}{j_1^{r_1}\cdots j_d^{r_d}}.
 \label{eq:MPL}
\end{equation}
The single-variable abbreviation
$\Li_{\mathbf r}(z)=\Li_{\mathbf r}(z,1,\ldots,1)$ is used below.
For $m\ge1$, set
\begin{equation}
 \Lcal_m(z)=\sum_{\mathbf r\vDash m}2^{\ell(\mathbf r)}\Li_{\mathbf r}(z),
 \qquad \Lcal_0(z)=1.
 \label{eq:Lm}
\end{equation}
The finite-harmonic generating function is
\begin{equation}
 \sum_{n\ge0}C_m(n)z^n=\frac{\Lcal_m(z)}{1-z},\qquad |z|<1.
 \label{eq:H-generating}
\end{equation}
Indeed, for each fixed decreasing tuple, the allowed upper bound $n$
runs from its largest member to infinity.

\begin{theorem}[Integral and colored-polylogarithm forms]
\label{thm:polylog}
For integers $m\ge0$, $p\ge1$, and $0<z<1$, put
\begin{equation}
 F_{m,p}(z)=\sum_{n\ge0}\frac{C_m(n)z^{2n+1}}{(2n+1)^p}.
 \label{eq:Fmp}
\end{equation}
Then
\begin{align}
 F_{m,p}(z)&=\frac1{(p-1)!}\int_0^z
        \log^{p-1}(z/t)\frac{\Lcal_m(t^2)}{1-t^2}\dd t,
        \label{eq:integral-master}\\
 F_{m,p}(z)&=2^{m-1}\sum_{\mathbf r\vDash m}
       \sum_{\epsilon_0,\ldots,\epsilon_d\in\{\pm1\}}
       \epsilon_0\Li_{p,r_1,\ldots,r_d}
                  (\epsilon_0z,\epsilon_1,\ldots,\epsilon_d),
        \label{eq:colored-master}
\end{align}
where $d=\ell(\mathbf r)$. Formula~\eqref{eq:colored-master} extends
to $|z|<1$ by analyticity. For $p\ge2$, both formulas extend to $z=1$
by ordinary absolute convergence. In particular, for $p=2k+1$, $k\ge1$,
Theorem~\ref{thm:harmonic} evaluates their endpoints in ordinary zeta
values and $\log2$ at every depth.
\end{theorem}
\begin{proof}
For each $n$,
\[
 \frac1{(p-1)!}\int_0^z t^{2n}\log^{p-1}(z/t)\dd t
       =\frac{z^{2n+1}}{(2n+1)^p}.
\]
Insert \eqref{eq:H-generating}. Positivity for real $z$, or absolute
convergence inside the unit disk, permits the interchange.

For the colored expression, fix a composition of length $d$. Replace
the inner harmonic indices $j_i$ by even integers $h_i=2j_i$ and
the outer denominator $2n+1$ by $h_0$. This contributes the factor
$2^m$ from the denominators. The inequalities become
$h_0>h_1>\cdots>h_d\ge1$, with $h_0$ odd and the others even.
The parity filters are
\[
 \mathbf1_{h_0\text{ odd}}=\frac12\sum_{\epsilon_0=\pm1}
                       \epsilon_0\epsilon_0^{h_0},\qquad
 \mathbf1_{h_i\text{ even}}=\frac12\sum_{\epsilon_i=\pm1}
                       \epsilon_i^{h_i}.
\]
Multiplication by the composition weight $2^d$ leaves
$2^m2^d/2^{d+1}=2^{m-1}$, proving \eqref{eq:colored-master}.
At $m=0$, the convention gives the familiar
$(\Li_p(z)-\Li_p(-z))/2$.

For fixed $m$, the coefficients grow at most logarithmically to
power $m$. Thus $p\ge2$ makes the endpoint sum absolutely convergent.
The corresponding multiple sums have first index at least two,
so every colored term is also absolutely convergent. This permits
separate endpoint limits; no regularization is being hidden here.
\end{proof}

\begin{corollary}[An exact dilogarithmic integral]
\label{cor:dilog-integral}
One has
\begin{align}
 &\frac12\int_0^1\frac{\log^2x}{1-x^2}
          [\Li_2(x^2)+\log^2(1-x^2)]\dd x\notag\\
 &\qquad=\frac{31}{8}\zeta(5)-\frac{45}{8}\log2\,\zeta(4)
       +\frac72\log^22\,\zeta(3)-\frac78\zeta(2)\zeta(3).
 \label{eq:dilog-integral}
\end{align}
\end{corollary}
\begin{proof}
The identity
\[
 \sum_{n\ge0}H_n^2z^n=
 \frac{\Li_2(z)+\log^2(1-z)}{1-z}
\]
follows by splitting the double harmonic sum into its diagonal and
two strict triangles (equivalently, $\Li_{1,1}(z)=\tfrac12\log^2(1-z)$).
Use the integral in \eqref{eq:integral-master} at $p=3$, then
\eqref{eq:example2}. The endpoint logarithms are integrable; the
preceding positive-sum argument also proves convergence directly.
\end{proof}

\subsection{Euler primitives and a Gamma primitive}
Set
\begin{equation}
 F_{m,0}(z)=\frac{z\Lcal_m(z^2)}{1-z^2}.
 \label{eq:Fm0}
\end{equation}
Termwise differentiation inside the disk gives the complete Euler
primitive ladder
\begin{equation}
 z\frac{\mathrm d}{\mathrm dz}F_{m,p}(z)=F_{m,p-1}(z),\qquad
 F_{m,p}(z)=\int_0^z\frac{F_{m,p-1}(t)}t\dd t.
 \label{eq:Euler-ladder}
\end{equation}
The basepoint zero fixes the integration constants. Formula
\eqref{eq:colored-master} supplies explicit primitives in multiple
polylogarithms, and the odd-denominator endpoint values for $p\ge3$
are supplied by \eqref{eq:all-harmonic}. The endpoint $p=1$ must not
be assigned a finite value by this argument.

At $u=0$, the general resolvent also has an elementary Gamma primitive:
\begin{equation}
 \sum_{n\ge0}\left[\arctanh\frac{v}{n+a}-\frac{v}{n+a}\right]
 =v\psi(a)-\frac12\log\frac{\Gamma(a+v)}{\Gamma(a-v)},
 \qquad |v|<a.
 \label{eq:gamma-primitive}
\end{equation}
Both sides vanish at zero. Their derivatives are the two sides of
\eqref{eq:resolvent} at $u=0$, by \eqref{eq:boundary0}.
The logarithm is the analytic branch normalized at $v=0$;
for real $|v|<a$ it is the real logarithm. Normal convergence follows
from the cubic cancellation in the summand.

\subsection{Stieltjes primitives with a fixed convention}
Hurwitz spectral derivatives provide a compatible primitive for each
boundary constant. Differentiating
$\partial_a\zeta(s,a)=-s\zeta(s+1,a)$ at $s=0$ using
\eqref{eq:stieltjes-convention} yields
\begin{equation}
 \partial_a\zeta^{(m+1)}(0,a)=(-1)^{m+1}(m+1)\gamma_m(a).
 \label{eq:stieltjes-derivative}
\end{equation}
Consequently, for $a,a_0>0$,
\begin{equation}
 \int_{a_0}^a\gamma_m(t)\dd t
 =\frac{(-1)^{m+1}}{m+1}
       [\zeta^{(m+1)}(0,a)-\zeta^{(m+1)}(0,a_0)].
 \label{eq:stieltjes-primitive}
\end{equation}
The superscript on $\zeta$ differentiates the spectral variable.
At $m=0$, Lerch's identity gives the expected
$-\log\Gamma(a)+\log\Gamma(a_0)$. These primitive identities are
classical consequences of Hurwitz-zeta calculus, included to tie the
boundary jets to the manuscript's integration conventions, not as a
claim of new priority.

\section{Zero-safe differentiation at terminating parameter values}
\label{sec:revival}
A family may terminate at a parameter value while its derivative has an
infinite tail. The normalized half-center makes this explicit without
special-function ambiguities.

\begin{proposition}[Revived derivatives of a terminating product]
\label{prop:revival}
For integers $M\ge1$ and $n\ge M$,
\begin{equation}
 P_n(M)=0,\qquad
 P_n'(M)=(-1)^M\frac{M!(M-1)!(n-M)!}{(n+M)!}.
 \label{eq:revival}
\end{equation}
For $0\le n<M$,
\begin{align}
 P_n(M)&=(-1)^n\frac{(M-1)!M!}{(M-1-n)!(M+n)!},\notag\\
 P_n'(M)&=P_n(M)[H_{M-1}-H_{M-1-n}-H_{M+n}+H_M].
 \label{eq:head-derivative}
\end{align}
\end{proposition}
\begin{proof}
For $n\ge M$, differentiate the unique zero factor $(M-u)/(M+u)$
in the finite product \eqref{eq:Pn}; all other factors can then be
evaluated at $M$. Their factorial products give \eqref{eq:revival}.
For $n<M$, no factor vanishes. Logarithmic differentiation of the
finite nonzero product gives \eqref{eq:head-derivative}.
No logarithm of a zero factor is taken.
\end{proof}

For $M=1$, all $n\ge1$ have $P_n'(1)=-1/(n(n+1))$.
One exact family appearing in the differentiated generator is
\begin{equation}
 \sum_{n\ge1}\frac1{n(n+1)(2n+1)^{2k+1}}
 =4k+3-4\log2
       -4\sum_{j=1}^k(1-2^{-2j-1})\zeta(2j+1),\quad k\ge1.
 \label{eq:revival-sum}
\end{equation}
It also has an independent elementary proof. With $x=2n+1$, repeatedly
use
\[
 \frac1{n(n+1)x^{2k+1}}
 =\frac1{n(n+1)x^{2k-1}}-\frac4{x^{2k+1}}.
\]
The initial sum at $k=0$ is $3-4\log2$, obtained from
$1/[n(n+1)(2n+1)]=1/n+1/(n+1)-4/(2n+1)$ by partial sums.
This proves \eqref{eq:revival-sum}.

The general warning about reciprocal-Gamma zeros was already explicit
in the centered continuation \cite{repo-centered}. The present formulas
are a concrete finite-product realization, not a claim that the warning
itself is new.

\section{Verification, audit conclusions, and integration}
\label{sec:verification}
\subsection{What the executable checks establish}
The package contains a finite exact coefficient engine and separate
high-precision diagnostics. The exact tests use rational residue classes
of Gamma shifts to construct divisors independently of the theorem's
parameter classification. They compare rationality and Laurent
termination on a grid of 10,356 valid parameter choices. The grid
uses $0\le N\le5$, quarter-integer positive centers through $3$, and
quarter-integer positive $b,c$ subject to $d_0>0$.

The full exact run executes 21,009 assertions: 10,356 rationality checks,
10,356 termination checks, 204 monomial-weight checks, four printed
harmonic-identity checks, 56 finite composition checks, 30 revived-zero
derivative checks, and three unweighted odd-denominator checks.
These finite calculations are regression tests of formulas and
implementations. They do not replace the all-parameter proofs.

The numerical checks compare independently summed series against Gamma,
polygamma, zeta, and Stieltjes expressions. They include an independent
quadrature of \eqref{eq:dilog-integral}. Series are evaluated by a finite
head plus an Euler--Maclaurin tail at 40 decimal digits. The assertion
threshold is a scaled residual of $10^{-28}$; it is not a rigorous error
bound. All 28 diagnostics passed; the largest observed scaled residual was
$5.511\times10^{-40}$. Full parameters, values, observed residuals,
software versions, and methods appear in the result files. No interval certification or
proof-assistant verification is claimed.

\subsection{Corrections and guards}
No false theorem was identified in the selected source ranges. The
following distinctions are necessary for correct integration.

\begin{enumerate}[leftmargin=*,itemsep=3pt]
\item The canonical first-resonance proof and the preserved all-resonance
proof have different editorial status. The present package does not
relabel the already proved all-resonance formula as a newly settled
conjecture.
\item A zero summed remainder is not finite termination.
Theorem~\ref{thm:zero-curves} proves this within the same family.
\item Formal termination needs an argument ruling out a flat remainder.
Lemma~\ref{lem:formal-exact} supplies that argument, rather than assuming
that a truncated asymptotic expansion equals the function.
\item Reciprocal-Gamma zeros are zeros of values, not automatically of
jets. Keep the entire factors or use finite products before differentiating.
\item The normalization constants are $\Gamma(u)P(u)/2$,
$[v^{2k}]Q_a(u,v)/(2u)$, and $[t^{2k}]\Gcal(u,t)/(4u)$.
The Stieltjes expansion uses $2u$, not $u$.
\item The all-depth theorem evaluates a weighted sum over compositions
and colors. It asserts neither a reduction of every individual colored
multiple zeta value nor a minimality or independence theorem.
\end{enumerate}

The suggested placement is an additive continuation of the centered
Dougall report, with the classification linked from the canonical
mixed-Dougall section. The harmonic/polylogarithmic sections can also be
cross-referenced from the depth and integration chapters. The supplied
integration note gives a compact insertion with distinct labels. Existing
source files should not be silently replaced by this draft.

\section{Further research questions}
\label{sec:further}
\paragraph{1. Interlaced resonances.}
Classify rationality and finite termination at $u=-N-1/2$ in the same
positive-parameter family. The shift multiset then has half-integral
total sum, so the fixed-class argument in
Theorem~\ref{thm:classification} changes. A full analysis should preserve
initial reciprocal-Gamma zeros and distinguish polynomial from Laurent
termination.

\paragraph{2. Larger Gamma quotients.}
For $x\prod_{i=1}^{2r}\Gamma(x+\alpha_i)/\Gamma(x+1-\alpha_i)$,
can rationality and Laurent termination be described by an effective
matching theorem and a finite arrangement of parameter loci? Modulo-integer
matching decides rationality, but a comparably short cancellation-free
classification may need additional structure beyond four factors.

\paragraph{3. Global geometry of exact zero remainders.}
The implicit curves in Theorem~\ref{thm:zero-curves} are local. Determine
the connected components and possible singularities of the zero set of
the collapsed digamma bracket inside the positive domain. Which components
meet the terminating locus, and which contain no rational points in the
Gamma-quotient sense?

\paragraph{4. Independent harmonic coordinates outside the Bell family.}
At each weight, which individual sums such as
$\sum H_n^3/(2n+1)^3$ or $\sum H_n^{(3)}/(2n+1)^3$ reduce to the
single-zeta/logarithm algebra? The present theorem evaluates a specific
combination and does not separate its components. A precise comparison
algebra is needed before a reduction or obstruction can be claimed.

\paragraph{5. General rational centers.}
For rational $a$, the finite formula \eqref{eq:finite-bell} lies in a
polygamma algebra at $a$. Transport this formula to a conductor-explicit
Dirichlet-$L$ and cyclotomic-polylogarithm basis, with zero-safe treatment
of principal-character poles. This is a reduction problem, not an
assertion of independence of the resulting coordinates.

\paragraph{6. Twisted and alternating resolvents.}
Determine which numerator twists $\omega^n$ admit a comparably short
Gamma evaluation, or a finite combination of such evaluations.
An alternating analogue would bear more directly on the manuscript's
Gaussian sums, but the nonalternating identities here do not imply it.

\paragraph{7. Relations to known multiple-zeta sum formulas.}
Give a complete comparison of \eqref{eq:colored-master} at $z=1$ with
published height, depth, and level-two sum formulas, including the
hypergeometric harmonic-differentiation literature. This is an important
priority and structural-comparison task left open by the limited source
review, not evidence that such antecedents are absent.

\paragraph{8. Beyond the first Euler endpoint.}
Can additional parameter symmetries evaluate the even-$p$ endpoints of
$F_{m,p}(1)$ in a comparably explicit algebra? The integral and colored
forms exist for every $p\ge2$, but the Gamma generator here singles out
odd $p\ge3$.

\paragraph{9. Formal verification of the finite interfaces.}
Formalize Gamma-shift divisor matching, the four-class parity argument,
the cancellation-free product classification, and the finite composition
identity. These finite interfaces are natural preliminary targets before
formalizing normal convergence and the Rogers--Dougall continuation.

\appendix
\section{A finite coefficient algorithm}
\label{app:algorithm}
Let $T=t^2$ and write
$\log\Gcal(u,t)=\sum_{r\ge1}L_r(T)u^r$.
For a requested $[u^{m+1}T^k]$ coefficient, discard powers of $T$ above
$k$ from every $L_r$, and retain only $r\le m+1$. Define
\begin{equation}
 G_0(T)=1,\qquad
 G_n(T)=\frac1n\sum_{r=1}^n rL_r(T)G_{n-r}(T)
                         \pmod{T^{k+1}}.
 \label{eq:coefficient-algorithm}
\end{equation}
Then $G_n(T)=[u^n]\Gcal(u,t)\pmod{T^{k+1}}$. The recurrence follows
by differentiating an exponential formal series and comparing coefficients.
It uses only finitely many rational operations. In particular,
\eqref{eq:all-harmonic} is $(-1)^m[T^k]G_{m+1}(T)/4$.

The executable implementation uses the formal symbols \texttt{L} for
$\log2$, \texttt{Zr} for $\zeta(r)$, and \texttt{T} for $t^2$.
The delivered coefficient table covers $0\le m\le5$, $1\le k\le3$.
Even zeta values are simplified only when explicitly requested, using
\begin{equation}
 \zeta(2j)=\frac{(-1)^{j+1}B_{2j}\,2^{2j}6^j}{2(2j)!}\zeta(2)^j.
 \label{eq:even-zeta}
\end{equation}
No integer-relation algorithm or floating-point fitting is used.

\section{Domains and conventions at a glance}
\begin{longtable}{@{}p{.29\linewidth}p{.65\linewidth}@{}}
\toprule
Identity or object & Domain and convention \\\midrule\endhead
Classification & $N\ge0$, $\kappa,b,c,d_0>0$, $d_0=1+2\kappa-b-c$.
Gamma denominators are entire reciprocal-Gamma factors.\\[4pt]
General resolvent & $a>0$, $|v|<a$, $-1<\Repart u<1$;
$u=0$ is removable.\\[4pt]
Half-center generator & $|t|<1$, $\Repart u>-1$;
Taylor coefficients are ordinary coefficients, not derivatives.\\[4pt]
Harmonic sum formula & Integers $m\ge0$, $k\ge1$; all sums are ordinary
and absolutely convergent.\\[4pt]
Polylogarithm endpoint & $p\ge2$; strictly decreasing nested indices.
The $p=1$ endpoint is not evaluated by this theorem.\\[4pt]
Stieltjes constants & $\zeta(1+s,a)=s^{-1}+\sum(-1)^m\gamma_m(a)s^m/m!$;
$\gamma_0(a)=-\psi(a)$.\\[4pt]
Gamma primitive & $|v|<a$, normalized analytic logarithm at $v=0$.\\[4pt]
Numerical evidence & 40-digit floating-point diagnostics; no certified
error intervals and no formal proof-assistant claim.\\
\bottomrule
\end{longtable}

\begin{thebibliography}{99}
\bibitem{dlmf16}
NIST Digital Library of Mathematical Functions, \S16.4,
\emph{Argument Unity}, especially Eq.~16.4.9, Rogers--Dougall summation.
\url{https://dlmf.nist.gov/16.4}. Accessed 11 October 2026.

\bibitem{dlmf5}
NIST Digital Library of Mathematical Functions, \S5.5,
\emph{Functional Relations}, and \S5.11, \emph{Asymptotic Expansions}.
\url{https://dlmf.nist.gov/5.5}; \url{https://dlmf.nist.gov/5.11}.
Accessed 11 October 2026.

\bibitem{dlmf25}
NIST Digital Library of Mathematical Functions, \S25.11,
\emph{Hurwitz Zeta Function}, and \S25.12, \emph{Polylogarithms}.
\url{https://dlmf.nist.gov/25.11}; \url{https://dlmf.nist.gov/25.12}.
Accessed 11 October 2026.

\bibitem{xu-wang}
Ce Xu and Weiping Wang, \emph{Dirichlet type extensions of Euler sums},
Comptes Rendus Math\'ematique \textbf{361} (2023), 979--1010.
DOI: \href{https://doi.org/10.5802/crmath.453}{10.5802/crmath.453}.
The publisher abstract and bibliographic information were inspected;
a full-paper priority comparison was not completed.

\bibitem{repo-main}
Vladimir Reshetnikov, \emph{ProveIt: Polylogarithms manuscript},
repository snapshot \pin.
\url{https://github.com/VladimirReshetnikov/ProveIt}.
Inspected files include \texttt{polylogarithms.tex} and selected chapters
under \texttt{Analysis/Polylogarithms/docs/manuscript/}.

\bibitem{repo-mixed}
ProveIt, \emph{Mixed Gamma derivatives and exact harmonic tail identities},
canonical manuscript file
\texttt{chapters/07-mixed-dougall-identities.tex}, lines 1--240,
at the snapshot in \cite{repo-main}.
The first-resonance and mixed-derivative results are antecedents here.

\bibitem{repo-centered}
ProveIt, \emph{Centered Dougall Resonance: Exact Stieltjes Jets,
Harmonic Identities, Polylogarithmic Subtraction, and Zeta Antiderivatives},
10 October 2026. Preserved continuation under
\texttt{Analysis/Polylogarithms/docs/reports/}\allowbreak
\texttt{stieltjes-harmonic-resonance/continuations/}\allowbreak
\texttt{centered-dougall-resonance/}.
Inspected: README; complete
\texttt{integration/centered-dougall-resonance.tex};
\texttt{article.tex}, lines 650--880, at the snapshot in \cite{repo-main}.
The all-resonance subtraction and collapse are attributed to this source.
\end{thebibliography}
\end{document}
